% Sugestões de texto para incluir na monografia
% Baseado nos resultados das análises geradas pelos scripts
% 
% Uso: Copie e adapte os trechos conforme necessário

%==============================================================================
% PARA A SEÇÃO: Fontes de Dados (Cap 2)
%==============================================================================

Os dados utilizados neste trabalho foram disponibilizados pelo projeto 
BikeScience, em parceria com a empresa Tembici, operadora do sistema de 
bicicletas compartilhadas de São Paulo. O banco de dados analisado contém 
\totalViagensDB{} viagens registradas entre \dataPrimeiraViagem{} e 
\dataUltimaViagem{}, abrangendo um período de \anosDataset{} anos.

Especificamente para o campus do Butantã da USP, foram identificadas 
\totalViagensUSP{} viagens envolvendo as \totalEstacoesUSP{} estações 
instaladas no campus. Destas, \viagensInternasUSP{} (\percentualInternas\%) 
correspondem a viagens internas ao campus (origem e destino dentro da USP), 
enquanto \viagensExternasUSP{} viagens têm uma das pontas fora do campus, 
evidenciando a integração do sistema de bicicletas compartilhadas com a 
mobilidade urbana mais ampla.

%==============================================================================
% PARA A SEÇÃO: Discussão - Diferenças entre Estações (Cap 5)
%==============================================================================

\subsection{Estações mais utilizadas}

A análise do volume de viagens por estação revela padrões claros de uso no 
campus. A Figura~\ref{fig:ranking-estacoes} apresenta o ranking completo das 
\totalEstacoesUSP{} estações USP ordenadas por volume total de viagens 
(saídas + chegadas).

\begin{figure}[htbp]
  \centering
  \includegraphics[width=0.95\textwidth]{ranking_estacoes_usp}
  \caption{Ranking de estações USP por volume de viagens (período: \dataPrimeiraViagem{} a \dataUltimaViagem{})}
  \label{fig:ranking-estacoes}
\end{figure}

As cinco estações mais movimentadas são:

\begin{enumerate}
  \item \textbf{\topEstacao1ID{} - \topEstacao1Nome{}}: \topEstacao1Viagens{} viagens
  \item \textbf{\topEstacao2ID{} - \topEstacao2Nome{}}: \topEstacao2Viagens{} viagens
  \item \textbf{\topEstacao3ID{} - \topEstacao3Nome{}}: \topEstacao3Viagens{} viagens
  \item \textbf{\topEstacao4ID{} - \topEstacao4Nome{}}: \topEstacao4Viagens{} viagens
  \item \textbf{\topEstacao5ID{} - \topEstacao5Nome{}}: \topEstacao5Viagens{} viagens
\end{enumerate}

A estação \topEstacao1Nome{} destaca-se como a mais utilizada, concentrando 
grande volume de viagens devido à sua localização central e proximidade com 
o restaurante universitário, ponto de convergência natural da comunidade 
acadêmica durante os horários de refeição.

A estação \topEstacao2Nome{} ocupa a segunda posição, atuando como principal 
ponto de acesso externo ao campus. Sua alta demanda reflete a integração do 
sistema de bicicletas compartilhadas com o transporte público metropolitano.

\subsection{Balanço entre saídas e chegadas}

A Figura~\ref{fig:balanco-estacoes} mostra o balanço entre saídas e chegadas 
para cada estação, onde valores positivos indicam mais partidas que chegadas 
(estações "geradoras") e valores negativos indicam o oposto (estações 
"atratoras").

\begin{figure}[htbp]
  \centering
  \includegraphics[width=0.95\textwidth]{balanco_estacoes_usp}
  \caption{Balanço de viagens por estação (verde = mais saídas, vermelho = mais chegadas)}
  \label{fig:balanco-estacoes}
\end{figure}

Observa-se que a estação \topEstacao2Nome{} apresenta balanço negativo 
significativo, funcionando como atratora, enquanto outras estações apresentam 
balanço mais equilibrado ou ligeiramente positivo.

%==============================================================================
% PARA A SEÇÃO: Discussão - Padrões Temporais (Cap 5)
%==============================================================================

\subsection{Distribuição por hora do dia}

A Figura~\ref{fig:viagens-hora} mostra a distribuição de viagens ao longo do 
dia. Observa-se claramente que o horário de pico ocorre às \horaPico{}, com 
\viagensHoraPico{} viagens registradas nesse período.

\begin{figure}[htbp]
  \centering
  \includegraphics[width=0.9\textwidth]{viagens_por_hora}
  \caption{Distribuição de viagens USP por hora do dia}
  \label{fig:viagens-hora}
\end{figure}

Este padrão sugere forte relação com os horários de refeições e intervalos 
entre aulas, evidenciando o uso do sistema como meio de locomoção durante o 
expediente acadêmico.

\subsection{Dias úteis vs fim de semana}

A análise temporal revela que \percentualDiasUteis\% das viagens ocorrem em 
dias úteis, com média de \mediaDiaUtil{} viagens por dia útil contra apenas 
\mediaFimSemana{} viagens por dia de fim de semana. 

\begin{figure}[htbp]
  \centering
  \includegraphics[width=0.9\textwidth]{viagens_por_dia_semana}
  \caption{Distribuição de viagens USP por dia da semana}
  \label{fig:viagens-dia-semana}
\end{figure}

Este padrão, ilustrado na Figura~\ref{fig:viagens-dia-semana}, confirma que 
o uso do sistema está fortemente vinculado às atividades acadêmicas regulares, 
com demanda significativamente reduzida aos finais de semana.

\subsection{Sazonalidade e período letivo}

A distribuição mensal de viagens (Figura~\ref{fig:viagens-mes}) evidencia 
clara sazonalidade, com \percentualPeriodoLetivo\% das viagens concentradas 
no período letivo. O mês de \mesPico{} apresenta o maior volume de viagens.

\begin{figure}[htbp]
  \centering
  \includegraphics[width=0.9\textwidth]{viagens_por_mes}
  \caption{Distribuição de viagens USP por mês (todos os anos agregados)}
  \label{fig:viagens-mes}
\end{figure}

Observa-se redução significativa em janeiro, fevereiro, julho e dezembro, 
correspondentes aos períodos de férias acadêmicas, reforçando a conclusão de 
que o sistema atende predominantemente à comunidade universitária em suas 
atividades regulares.

%==============================================================================
% PARA A SEÇÃO: Discussão - Fluxos (Cap 5)
%==============================================================================

\subsection{Rotas mais frequentes}

A análise de fluxos origem-destino identifica os principais pares de estações 
conectados por viagens. A Figura~\ref{fig:top-rotas} apresenta as 20 rotas 
mais frequentes.

\begin{figure}[htbp]
  \centering
  \includegraphics[width=0.95\textwidth]{top20_rotas_frequentes}
  \caption{Top 20 rotas mais frequentes entre estações}
  \label{fig:top-rotas}
\end{figure}

As três rotas mais utilizadas são:

\begin{enumerate}
  \item \textbf{\topRota1OrigemNome{} → \topRota1DestinoNome{}}: \topRota1Viagens{} viagens
  \item \textbf{\topRota2OrigemNome{} → \topRota2DestinoNome{}}: \topRota2Viagens{} viagens
  \item \textbf{\topRota3OrigemNome{} → \topRota3DestinoNome{}}: \topRota3Viagens{} viagens
\end{enumerate}

Estas rotas revelam os principais eixos de deslocamento no campus, conectando 
pontos de alimentação, estudo e acesso externo.

\subsection{Fluxos internos vs externos}

Do total de \totalViagensUSP{} viagens envolvendo o campus, \viagensInternasUSP{} 
(\percentualInternas\%) são viagens internas (ambas as estações dentro da USP), 
enquanto \viagensExternasUSP{} viagens conectam o campus com a cidade.

\begin{figure}[htbp]
  \centering
  \includegraphics[width=0.9\textwidth]{fluxos_internos_externos}
  \caption{Distribuição entre viagens internas e externas ao campus}
  \label{fig:fluxos-internos-externos}
\end{figure}

A predominância de viagens internas (Figura~\ref{fig:fluxos-internos-externos}) 
demonstra que o sistema de bicicletas compartilhadas desempenha papel importante 
na mobilidade interna do campus, facilitando deslocamentos entre diferentes 
unidades acadêmicas e pontos de serviço.

%==============================================================================
% PARA A SEÇÃO: Discussão - Impacto COVID-19 (Cap 5)
%==============================================================================

\subsection{Impacto da pandemia COVID-19}

A Figura~\ref{fig:evolucao-temporal} apresenta a evolução mensal das viagens 
ao longo do período analisado, evidenciando o impacto da pandemia de COVID-19.

\begin{figure}[htbp]
  \centering
  \includegraphics[width=0.95\textwidth]{evolucao_temporal_viagens}
  \caption{Evolução temporal das viagens USP (2020-2022), destacando o período de pandemia}
  \label{fig:evolucao-temporal}
\end{figure}

Observa-se queda abrupta no volume de viagens a partir de março de 2020, 
coincidindo com as medidas de isolamento social e suspensão das atividades 
presenciais. A retomada gradual é visível ao longo de 2021 e 2022, acompanhando 
o retorno progressivo das atividades acadêmicas presenciais.

%==============================================================================
% PARA A SEÇÃO: Conclusão
%==============================================================================

Este trabalho analisou \totalViagensUSP{} viagens de bicicletas compartilhadas 
envolvendo o campus do Butantã da USP, identificando padrões claros de uso:

\begin{itemize}
  \item \textbf{Uso predominantemente acadêmico}: \percentualDiasUteis\% das 
        viagens ocorrem em dias úteis e \percentualPeriodoLetivo\% em período letivo
  
  \item \textbf{Mobilidade interna}: \percentualInternas\% das viagens são 
        internas ao campus, evidenciando o papel do sistema na conexão entre 
        diferentes pontos do campus
  
  \item \textbf{Padrões horários}: Pico de uso às \horaPico{}, coincidindo 
        com horários de intervalo e refeições
  
  \item \textbf{Concentração em poucas estações}: A estação \topEstacao1Nome{} 
        concentra \topEstacao1Viagens{} viagens, atuando como hub principal
\end{itemize}

%==============================================================================
% FIM DAS SUGESTÕES
%==============================================================================
